\documentclass[11pt,a4paper]{article}
\usepackage[T1]{fontenc}
\usepackage[utf8]{inputenc}
\usepackage{lmodern}
\usepackage{microtype}
\usepackage{geometry}
\usepackage{listings}
\usepackage{xcolor}
\usepackage{hyperref}
\usepackage{parskip}
\usepackage{booktabs}

\geometry{a4paper, margin=2.5cm}

\definecolor{codebg}{RGB}{245,245,245}
\definecolor{codeframe}{RGB}{200,200,200}
\definecolor{keyword}{RGB}{0,100,180}
\definecolor{comment}{RGB}{100,140,100}
\definecolor{string}{RGB}{180,60,60}

\lstset{
  backgroundcolor=\color{codebg},
  frame=single,
  rulecolor=\color{codeframe},
  basicstyle=\ttfamily\small,
  keywordstyle=\color{keyword}\bfseries,
  commentstyle=\color{comment}\itshape,
  stringstyle=\color{string},
  breaklines=true,
  numbers=left,
  numberstyle=\tiny\color{gray},
  numbersep=6pt,
  showstringspaces=false,
  tabsize=4,
}

\hypersetup{
  colorlinks=true,
  linkcolor=blue,
  urlcolor=blue,
}

\title{\textbf{Wilo-Pumpe Teil 3: PIO misst alles}\\[4pt]
  \large Hardware statt Software --- wie der RP2040-PIO HIGH- und LOW-Zeit\\
  direkt in Hardware misst und warum 250\,kHz dabei Routine sind}
\author{make Magazin $\cdot$ Projekt-Doku}
\date{April 2026 \quad \texttt{github.com/gerontec/wilo\_pwm}}

\begin{document}

\maketitle
\thispagestyle{empty}

\noindent\rule{\textwidth}{0.4pt}
\begin{center}
\textit{Teil 1 nutzte C-Interrupts, Teil 2 Core~1-Polling.
Teil 3 delegiert die gesamte Messung an den Hardware-PIO ---
kein CPU-Overhead, kein Debounce, keine Firmware-Aenderung noetig.
Ergebnis: stabile Pulsmessung bis 250\,kHz als reines Python-Modul.}
\end{center}
\noindent\rule{\textwidth}{0.4pt}

\bigskip

% -----------------------------------------------------------------------
\section{Warum PIO besser als Core-1-Polling ist}

In Teil\,2 lief eine enge Polling-Schleife auf Core~1 mit
26\,Millionen Iterationen pro Sekunde.
Das funktioniert, hat aber zwei Schwaechen:

\begin{itemize}
\item \textbf{Software-Debounce noetig.}
  Jeder Glitch auf der Leitung muss manuell gefiltert werden ---
  mit adaptivem Schwellwert, der bei Frequenzaenderungen neu berechnet wird.
\item \textbf{Firmware-Build erforderlich.}
  Das Core-1-Modul ist als \texttt{USER\_C\_MODULES} in die Firmware gelinkt ---
  jede Aenderung erfordert einen neuen Build und BOOTSEL-Flash.
\end{itemize}

Der RP2040-PIO loest beide Probleme auf einmal.
PIO (Programmable I/O) ist ein eigenstaendiger Prozessor mit
eigenem Instruktionssatz und eigenem Takt, voellig unabhaengig
von ARM-Core~0 und Core~1.
Ein PIO-Programm laeuft deterministisch Zyklus fuer Zyklus,
wird nicht von GC-Pausen unterbrochen und braucht keine IRQ-Vektoren.
Und das Beste: In MicroPython laesst sich PIO als reines Python-Modul
deployen --- kein Compiler, kein Flash.

% -----------------------------------------------------------------------
\section{Das PIO-Assembler-Programm}

Das Messprinzip ist direkt:
\begin{enumerate}
  \item Auf steigende Flanke warten (\texttt{wait})
  \item Abwaertszaehler X von \texttt{0xFFFFFFFF} herunterzaehlen solange
        der Pin HIGH ist
  \item Zaehlerstand in RX-FIFO schieben (\texttt{push})
  \item Dasselbe fuer die LOW-Phase wiederholen
\end{enumerate}

\begin{lstlisting}[language=Python, caption={pwmfeedback\_pio.py --- PIO-Programm fuer HIGH/LOW-Messung}]
@rp2.asm_pio()
def _pwm_measure():
    pull()           # Y-Initialwert (0xFFFFFFFF) aus TX-FIFO laden
    mov(y, osr)

    wrap_target()

    wait(0, pin, 0)  # Sync: warten bis Pin LOW
    wait(1, pin, 0)  # warten auf steigende Flanke (Pin wird HIGH)

    # HIGH-Zeit messen: X abwaerts zaehlen solange Pin HIGH
    mov(x, y)                    # X = 0xFFFFFFFF
    label("high_loop")
    jmp(pin, "high_count")       # Pin HIGH? -> weiterzaehlen
    jmp("high_done")             # Pin LOW?  -> fertig
    label("high_count")
    jmp(x_dec, "high_loop")      # X--, und zurueck
    label("high_done")
    mov(isr, x)
    push(noblock)                # Zaehlerstand -> RX-FIFO

    # LOW-Zeit messen: X abwaerts zaehlen solange Pin LOW
    mov(x, y)
    label("low_loop")
    jmp(pin, "low_done")         # Pin HIGH? -> LOW-Phase vorbei
    jmp(x_dec, "low_loop")       # X--, weiterzaehlen
    label("low_done")
    mov(isr, x)
    push(noblock)

    wrap()
\end{lstlisting}

Das Programm besteht aus 15 Instruktionen --- weit unter dem Limit von 32.
Die Zaehlschleife benoetigt \textbf{2 Zyklen pro Schritt}.
Bei einem PIO-Takt von 1\,MHz ergibt das eine Zeitaufloesung von 2\,\textmu s.

% -----------------------------------------------------------------------
\section{State Machine und Python-Auswertung}

Die State Machine wird mit einem einzigen \texttt{put()} initialisiert,
der den Startwert 0xFFFFFFFF in den TX-FIFO legt:

\begin{lstlisting}[language=Python, caption={Initialisierung der State Machine}]
pin = Pin(5, Pin.IN, Pin.PULL_UP)
sm = rp2.StateMachine(
    0, _pwm_measure,
    freq=1_000_000,    # 1 MHz -> 1 us/Zyklus, 2 us/Zaehlschritt
    in_base=pin,
    jmp_pin=pin,
)
sm.put(0xFFFFFFFF)     # Y-Initialwert
sm.active(1)
\end{lstlisting}

Die Python-Seite liest HIGH- und LOW-Wert paarweise aus dem FIFO
und rechnet Zeitdauer, Frequenz und Duty-Cycle:

\begin{lstlisting}[language=Python, caption={Auswertung: Frequenz und Duty-Cycle aus Zaehlerstand}]
while sm.rx_fifo() >= 2:
    x_high = sm.get()
    x_low  = sm.get()
    # (0xFFFFFFFF - x) * 2 us/Count = Pulsdauer in Mikrosekunden
    high_us = (0xFFFFFFFF - x_high) * 2
    low_us  = (0xFFFFFFFF - x_low)  * 2
    T       = high_us + low_us
    freq    = 1_000_000.0 / T        # Hz
    duty    = high_us / T * 100.0    # Prozent
\end{lstlisting}

\texttt{push(noblock)} sorgt dafuer, dass veraltete Werte verworfen werden,
wenn der FIFO voll ist --- Python kann die Daten im eigenen Rhythmus
(alle 5\,s) abholen, ohne dass die State Machine blockiert.

% -----------------------------------------------------------------------
\section{Frequenzlimit: Wann stoesst PIO an seine Grenzen?}

Die Zaehlschleife benoetigt 2 PIO-Zyklen pro Schritt.
Bei einem PIO-Takt von 1\,MHz gilt:

\[
  t_{\min} = 2 \times \frac{1}{f_{\text{PIO}}} = 2\,\mu\text{s}
\]

Fuer ein Signal mit 50\,\% Duty-Cycle muss sowohl HIGH als auch LOW
mindestens 2\,\textmu s betragen --- also Periode $\geq$ 4\,\textmu s,
Frequenz $\leq$ \textbf{250\,kHz}.

Bei niedrigem Duty-Cycle (z.\,B.\ 5\,\%) ist der HIGH-Impuls kurz:
$0{,}05 \times T$.
Das Minimum von 2\,\textmu s HIGH-Dauer limitiert dann die Frequenz:
$f_{\max} = 1 / (2\,\mu\text{s} / 0{,}05) = \mathbf{25\,kHz}$.

Wer hoehere Frequenzen braucht, erhoet den PIO-Takt auf bis zu 125\,MHz
(Systemtakt):

\begin{lstlisting}[language=Python]
sm = rp2.StateMachine(0, _pwm_measure, freq=125_000_000, ...)
# -> 16 ns/Zyklus, 32 ns/Zaehlschritt -> max ~31 MHz bei 50% Duty
\end{lstlisting}

\begin{center}
\begin{tabular}{lll}
\toprule
\textbf{PIO-Takt} & \textbf{Aufloesung} & \textbf{Max. Freq. (50\% Duty)} \\
\midrule
1\,MHz   & 2\,\textmu s & 250\,kHz \\
10\,MHz  & 200\,ns      & 2{,}5\,MHz \\
125\,MHz & 16\,ns       & $\sim$31\,MHz \\
\bottomrule
\end{tabular}
\end{center}

Fuer 75\,Hz (Wilo-Pumpe) ist der 1-MHz-Takt 3000-fach ueberdimensioniert ---
die Messung ist absolut stabil.

% -----------------------------------------------------------------------
\section{Vergleich aller vier Methoden}

\begin{center}
\begin{tabular}{lllll}
\toprule
\textbf{Methode} & \textbf{Max. Freq.} & \textbf{Debounce} & \textbf{CPU-Last} & \textbf{Build} \\
\midrule
Python IRQ (pure)    & $\sim$2\,kHz    & noetig & hoch      & nein \\
C-Natmod IRQ         & $\sim$10\,kHz   & noetig & mittel    & .mpy \\
Core~1 Polling       & $\sim$250\,kHz  & adaptiv & Core~1  & Firmware \\
PIO @ 1\,MHz         & 250\,kHz        & nicht noetig & null & nein \\
PIO @ 125\,MHz       & $\sim$31\,MHz   & nicht noetig & null & nein \\
\bottomrule
\end{tabular}
\end{center}

PIO ist in allen Kategorien besser als die Software-Ansaetze ---
und benoetigt als einzige Methode weder einen Firmware-Build
noch einen BOOTSEL-Flash.

% -----------------------------------------------------------------------
\section{Drop-in-Ersatz ohne Aenderung an main.py}

Das Modul \texttt{pwmfeedback\_pio.py} implementiert dieselbe API
wie alle Vorgaenger:

\begin{lstlisting}[language=Python, caption={Umschalten in main.py: eine Zeile}]
# vorher:
import pwmfeedback_core1 as pwmfeedback   # Core 1 Polling
# nachher:
import pwmfeedback_pio as pwmfeedback     # PIO Hardware
\end{lstlisting}

\texttt{init\_feedback\_pin()} startet die State Machine,
\texttt{get\_pump\_feedback()} liefert dasselbe Dict mit
\texttt{PIN5\_Freq\_Hz}, \texttt{PumpDuty} und \texttt{PumpStatus}.
Der Rest von \texttt{main.py} --- MQTT, Watchdog, Rampe, Boost ---
bleibt unveraendert.

Deploy per WebREPL oder mpremote, kein BOOTSEL noetig:

\begin{lstlisting}[language=bash, caption={Deploy ohne Firmware-Flash}]
mpremote connect /dev/ttyACM0 \
    cp pwmfeedback_pio.py :pwmfeedback_pio.py + \
    cp main.py :main.py + reset
# -> sofort aktiv, keine neue Firmware
\end{lstlisting}

% -----------------------------------------------------------------------
\section{Messwerte im Betrieb}

Gemessene Werte der Wilo iPWM-Pumpe im Normalbetrieb:

\begin{center}
\begin{tabular}{ll}
\toprule
\textbf{Kenngroesse} & \textbf{Wert} \\
\midrule
Frequenz    & 75{,}0\,Hz (stabil, keine Transienten) \\
Duty-Cycle  & 29{,}5--42\,\% (lastabhaengig) \\
Aufloesung  & 2\,\textmu s \\
CPU-Last    & 0\,\% (reine Hardware) \\
\bottomrule
\end{tabular}
\end{center}

Die Glitch-Probleme aus Teil\,2 --- kurze Transienten bei
153\,Hz oder 181\,Hz durch falsche Debounce-Timing ---
treten mit dem PIO-Ansatz nicht auf.
Die \texttt{wait()}-Instruktion synchronisiert auf die echte Flanke,
der Zaehler laeuft deterministisch ohne Unterbrechung.

\bigskip
\noindent\rule{\textwidth}{0.4pt}
\textbf{Code, Firmware \& Deploy-Script:}
\url{https://github.com/gerontec/wilo_pwm}

\end{document}
